% ============================================================================
% SOLUCIONES DEL MÓDULO 8. REPORTES AUTOMÁTICOS CON R MARKDOWN
% ============================================================================

\begin{solucion}{m8-1}
El bloque que lee los datos lleva \codigo{include = FALSE} para que no se
vea; la fecha es código en línea en el YAML, y la frase usa código en línea
con \codigo{nrow()} y \codigo{sum()} sobre la columna lógica
\codigo{es\_premium} (cada \codigo{TRUE} cuenta como 1). Como el archivo se
guarda en \archivo{reportes/}, los datos se buscan también en
\archivo{../datasets}.

\begin{textocodigo}[title=reportes/ej1-clientes.Rmd]
---
title: "Reporte de clientes"
author: "Tu nombre"
date: "`r format(Sys.Date(), '%d/%m/%Y')`"
output: html_document
---

```{r datos, include = FALSE}
library(dplyr)
library(readr)
library(knitr)
carpeta <- if (dir.exists("datasets")) "datasets" else "../datasets"
clientes <- read_csv(file.path(carpeta, "clientes.csv"),
                     show_col_types = FALSE)
```

## Resumen

La base tiene **`r nrow(clientes)` clientes**, de los cuales
**`r sum(clientes$es_premium)`** son premium
(`r round(mean(clientes$es_premium) * 100, 1)` %).

## Clientes por ciudad

```{r tabla-ciudades, echo = FALSE}
clientes %>%
  count(ciudad, sort = TRUE) %>%
  kable(col.names = c("Ciudad", "Clientes"),
        caption = "Clientes por ciudad")
```
\end{textocodigo}

Lo generamos con \menu{Knit} (o con
\codigo{rmarkdown::render("reportes/ej1-clientes.Rmd")}). Para comprobar
qué dirá el reporte, ejecutamos los mismos cálculos en la consola:

\begin{rcode}
library(tidyverse)
library(knitr)
clientes <- read_csv("datasets/clientes.csv", show_col_types = FALSE)
\end{rcode}

\begin{rcode}
nrow(clientes)
sum(clientes$es_premium)
round(mean(clientes$es_premium) * 100, 1)
clientes %>%
  count(ciudad, sort = TRUE) %>%
  kable(format = "pipe", col.names = c("Ciudad", "Clientes"))
\end{rcode}
\begin{rsalida}
[1] 200
[1] 73
[1] 36.5


|Ciudad      | Clientes|
|:-----------|--------:|
|Puebla      |       28|
|CDMX        |       27|
|Guadalajara |       22|
|Cancún      |       20|
|Toluca      |       20|
|León        |       18|
|Mérida      |       17|
|Querétaro   |       17|
|Tijuana     |       17|
|Monterrey   |       14|
\end{rsalida}

El reporte generado dice que hay 200 clientes, 73 premium
(36.5\,\%), y la tabla empieza con Puebla (28 clientes) y CDMX (27).
\end{solucion}

\begin{solucion}{m8-2}
Unimos las transacciones con la categoría del producto, resumimos, ordenamos
\emph{antes} de formatear (porque al formatear las columnas se vuelven texto)
y damos formato de negocio con \paquete{scales}.

\begin{rcode}
library(tidyverse)
library(knitr)
library(scales)

transacciones <- read_csv("datasets/transacciones.csv",
                          show_col_types = FALSE)
productos <- read_csv("datasets/productos.csv", show_col_types = FALSE)

por_categoria <- transacciones %>%
  left_join(productos %>% select(producto_id, categoria),
            by = "producto_id") %>%
  group_by(categoria) %>%
  summarise(ventas = sum(total_transaccion),
            unidades = sum(cantidad), .groups = "drop") %>%
  arrange(desc(ventas)) %>%                  # ordenar con números
  mutate(participacion = percent(ventas / sum(ventas), accuracy = 0.1),
         ventas = dollar(ventas, accuracy = 1),
         unidades = comma(unidades))

kable(por_categoria, format = "pipe",
      col.names = c("Categoría", "Ventas", "Unidades", "% del total"),
      align = c("l", "r", "r", "r"),
      caption = "Ventas 2023 por categoría")
\end{rcode}
\begin{rsalida}

Table: Ventas 2023 por categoría

|Categoría   |   Ventas| Unidades| % del total|
|:-----------|--------:|--------:|-----------:|
|Belleza     | $182,143|      551|       19.0%|
|Electrónica | $132,672|      403|       13.8%|
|Juguetes    | $125,093|      356|       13.0%|
|Alimentos   | $121,543|      356|       12.7%|
|Automotriz  | $104,085|      325|       10.8%|
|Deportes    |  $86,171|      280|        9.0%|
|Hogar       |  $63,407|      187|        6.6%|
|Mascotas    |  $57,668|      201|        6.0%|
|Libros      |  $57,162|      204|        6.0%|
|Ropa        |  $30,048|      106|        3.1%|
\end{rsalida}

Belleza es la categoría que más vende: \$182,143, el 19.0\,\% del total, seguida
de Electrónica (13.8\,\%). Ropa queda al final con solo el 3.1\,\%. Observa
la alineación: los números a la derecha permiten comparar magnitudes de un
vistazo.
\end{solucion}

\begin{solucion}{m8-3}
La función calcula primero el cumplimiento (\codigo{ventas / meta}) y con
\codigo{if}/\codigo{else if}/\codigo{else} elige la frase. Para aplicarla a
cada trimestre usamos \codigo{map\_chr()}, que llama a la función con cada
valor de ventas.

\begin{rcode}
library(tidyverse)
library(scales)

frase_meta <- function(ventas, meta) {
  cumplimiento <- ventas / meta
  montos <- paste0("(ventas de ", dollar(ventas, accuracy = 1),
                   " contra meta de ", dollar(meta, accuracy = 1), ").")
  if (ventas >= meta) {
    paste("Se superó la meta en",
          percent(cumplimiento - 1, accuracy = 0.1), montos)
  } else if (cumplimiento >= 0.9) {
    paste("Se alcanzó el", percent(cumplimiento, accuracy = 0.1),
          "de la meta", montos)
  } else {
    paste("Alerta: solo se alcanzó el",
          percent(cumplimiento, accuracy = 0.1), "de la meta", montos)
  }
}

# Los tres casos
frase_meta(110000, 100000)
frase_meta(95000, 100000)
frase_meta(70000, 100000)

# Aplicación a los trimestres de 2023
transacciones <- read_csv("datasets/transacciones.csv",
                          show_col_types = FALSE)
trimestres <- transacciones %>%
  group_by(trimestre) %>%
  summarise(ventas = sum(total_transaccion), .groups = "drop") %>%
  mutate(frase = map_chr(ventas, frase_meta, meta = 220000))
cat(paste0(trimestres$trimestre, ": ", trimestres$frase), sep = "\n")
\end{rcode}
\begin{rsalida}
[1] "Se superó la meta en 10.0% (ventas de $110,000 contra meta de $100,000)."
[1] "Se alcanzó el 95.0% de la meta (ventas de $95,000 contra meta de $100,000)."
[1] "Alerta: solo se alcanzó el 70.0% de la meta (ventas de $70,000 contra meta de $100,000)."
Q1: Se superó la meta en 8.8% (ventas de $239,390 contra meta de $220,000).
Q2: Se superó la meta en 8.8% (ventas de $239,457 contra meta de $220,000).
Q3: Se superó la meta en 6.5% (ventas de $234,241 contra meta de $220,000).
Q4: Se superó la meta en 12.2% (ventas de $246,905 contra meta de $220,000).
\end{rsalida}

Los cuatro trimestres superaron la meta de \$220,000; el mejor fue Q4, con
12.2\,\% por encima. Si cambias la meta a \$240,000, Q1, Q2 y Q3 pasan a
``Se alcanzó el \ldots'' sin tocar el texto: la frase depende solo de los
datos.
\end{solucion}

\begin{solucion}{m8-4}
El bloque \codigo{configuracion} fija las opciones globales; el bloque de
cálculo lleva \codigo{include = FALSE}; la gráfica define su tamaño y su
pie; y el bloque final usa \codigo{results = "asis"} para que el texto de
\codigo{cat()} se convierta en una lista.

\begin{textocodigo}[title=reportes/ej4-opciones.Rmd]
---
title: "Ventas mensuales 2023"
output:
  html_document: default
  word_document: default
---

```{r configuracion, include = FALSE}
knitr::opts_chunk$set(echo = FALSE, message = FALSE, warning = FALSE)
library(dplyr)
library(readr)
library(ggplot2)
library(scales)
```

```{r calculos, include = FALSE}
carpeta <- if (dir.exists("datasets")) "datasets" else "../datasets"
ventas <- read_csv(file.path(carpeta, "ventas_retail.csv"),
                   show_col_types = FALSE)
por_mes <- ventas %>%
  group_by(mes) %>%
  summarise(ventas = sum(total), .groups = "drop")
```

## Tendencia

```{r grafica, fig.cap = "Ventas por mes, 2023", fig.width = 7,
    fig.height = 3, out.width = "80%"}
ggplot(por_mes, aes(x = mes, y = ventas, group = 1)) +
  geom_line(color = "#2a78d6", linewidth = 0.8) +
  scale_y_continuous(labels = dollar) +
  labs(x = NULL, y = "Ventas") +
  theme_minimal() +
  theme(panel.grid.minor = element_blank(),
        axis.text.x = element_text(angle = 45, hjust = 1))
```

## Los tres mejores meses

```{r top-meses, results = "asis"}
top3 <- por_mes %>% arrange(desc(ventas)) %>% slice_head(n = 3)
cat(paste0("- **", top3$mes, "**: ", dollar(top3$ventas, accuracy = 1)),
    sep = "\n")
```
\end{textocodigo}

Se genera en los dos formatos con:

\begin{rnoejecutar}
rmarkdown::render("reportes/ej4-opciones.Rmd", output_format = "all")
\end{rnoejecutar}

La lista que escribe el último bloque es exactamente la que produce este
código en la consola:

\begin{rcode}
library(tidyverse)
library(scales)
ventas <- read_csv("datasets/ventas_retail.csv", show_col_types = FALSE)
por_mes <- ventas %>%
  group_by(mes) %>%
  summarise(ventas = sum(total), .groups = "drop")
top3 <- por_mes %>% arrange(desc(ventas)) %>% slice_head(n = 3)
cat(paste0("- **", top3$mes, "**: ", dollar(top3$ventas, accuracy = 1)),
    sep = "\n")
\end{rcode}
\begin{rsalida}
- **2023-07**: $54,056
- **2023-11**: $50,398
- **2023-05**: $48,108
\end{rsalida}

En el HTML y en el Word esas líneas aparecen como una lista con viñetas y
los meses en negritas, sin recuadro gris ni \codigo{\#\#}.
\end{solucion}

\begin{solucion}{m8-5}
El parámetro se declara en el YAML y se usa como \codigo{params\$categoria}
en el filtro, en el título y en las frases.

\begin{textocodigo}[title=reportes/ej5-categoria.Rmd]
---
title: "Reporte de la categoría `r params$categoria`"
output: html_document
params:
  categoria: "Electrónica"
---

```{r configuracion, include = FALSE}
knitr::opts_chunk$set(echo = FALSE, message = FALSE, warning = FALSE)
library(dplyr)
library(readr)
library(ggplot2)
library(scales)
library(knitr)
carpeta <- if (dir.exists("datasets")) "datasets" else "../datasets"
transacciones <- read_csv(file.path(carpeta, "transacciones.csv"),
                          show_col_types = FALSE)
productos <- read_csv(file.path(carpeta, "productos.csv"),
                      show_col_types = FALSE)
ventas_cat <- transacciones %>%
  left_join(productos, by = "producto_id") %>%
  filter(categoria == params$categoria)
if (nrow(ventas_cat) == 0) stop("Sin ventas de ", params$categoria)
```

En 2023, la categoría **`r params$categoria`** vendió
**`r dollar(sum(ventas_cat$total_transaccion), accuracy = 1)`** en
`r nrow(ventas_cat)` transacciones.

```{r top-productos}
ventas_cat %>%
  group_by(nombre_producto) %>%
  summarise(unidades = sum(cantidad),
            ventas = sum(total_transaccion), .groups = "drop") %>%
  arrange(desc(ventas)) %>%
  slice_head(n = 5) %>%
  mutate(ventas = dollar(ventas, accuracy = 1)) %>%
  kable(col.names = c("Producto", "Unidades", "Ventas"),
        align = c("l", "r", "r"),
        caption = "Productos más vendidos")
```

```{r grafica, fig.cap = "Ventas por mes", fig.height = 3}
ventas_cat %>%
  group_by(mes) %>%
  summarise(ventas = sum(total_transaccion), .groups = "drop") %>%
  ggplot(aes(x = mes, y = ventas)) +
  geom_col(fill = "#2a78d6", width = 0.7) +
  scale_y_continuous(labels = dollar) +
  labs(x = NULL, y = "Ventas") +
  theme_minimal() +
  theme(panel.grid.minor = element_blank(),
        axis.text.x = element_text(angle = 45, hjust = 1))
```
\end{textocodigo}

Para generar dos versiones, se recorre un vector de categorías y se da un
\codigo{output\_file} distinto a cada una (si no, la segunda sobrescribe a
la primera):

\begin{rnoejecutar}
for (cat_actual in c("Ropa", "Hogar")) {
  rmarkdown::render(
    "reportes/ej5-categoria.Rmd",
    params        = list(categoria = cat_actual),
    output_file   = paste0("categoria_", tolower(cat_actual), ".html"),
    output_dir    = "resultados",
    knit_root_dir = getwd(),
    envir         = new.env()
  )
}
\end{rnoejecutar}

Las frases de cada reporte coinciden con este cálculo:

\begin{rcode}
library(tidyverse)
library(scales)
transacciones <- read_csv("datasets/transacciones.csv",
                          show_col_types = FALSE)
productos <- read_csv("datasets/productos.csv", show_col_types = FALSE)
transacciones %>%
  left_join(productos, by = "producto_id") %>%
  filter(categoria %in% c("Ropa", "Hogar")) %>%
  group_by(categoria) %>%
  summarise(ventas = dollar(sum(total_transaccion), accuracy = 1),
            transacciones = n(), .groups = "drop")
\end{rcode}
\begin{rsalida}
# A tibble: 2 x 3
  categoria ventas  transacciones
  <chr>     <chr>           <int>
1 Hogar     $63,407            64
2 Ropa      $30,048            38
\end{rsalida}

Se obtienen \archivo{resultados/categoria\_ropa.html} y
\archivo{resultados/categoria\_hogar.html}, cada uno con su título, su frase,
su tabla y su gráfica.
\end{solucion}

\begin{solucion}{m8-6}
\codigo{expand\_grid()} crea las 40 combinaciones; con \codigo{count()}
contamos las transacciones por tienda y trimestre y las unimos con
\codigo{left\_join()}. Si alguna combinación no tuviera transacciones,
quedaría \codigo{NA}, que convertimos en 0 con \codigo{replace\_na()}.

\begin{rcode}
library(tidyverse)
transacciones <- read_csv("datasets/transacciones.csv",
                          show_col_types = FALSE)
tiendas <- read_csv("datasets/tiendas.csv", show_col_types = FALSE)

limpiar_nombre <- function(x) {
  x <- chartr("áéíóúüñÁÉÍÓÚÜÑ", "aeiouunAEIOUUN", x)
  x <- gsub("[^A-Za-z0-9]+", "-", x)
  tolower(gsub("^-|-$", "", x))
}

conteo <- transacciones %>% count(tienda_id, trimestre,
                                  name = "n_transacciones")

trabajos <- expand_grid(tienda_id = tiendas$tienda_id,
                        trimestre = paste0("Q", 1:4)) %>%
  left_join(tiendas %>% select(tienda_id, nombre_tienda),
            by = "tienda_id") %>%
  left_join(conteo, by = c("tienda_id", "trimestre")) %>%
  mutate(n_transacciones = replace_na(n_transacciones, 0),
         archivo = sprintf("2023-%s_T%02d_%s.pdf", trimestre, tienda_id,
                           limpiar_nombre(nombre_tienda)),
         generar = n_transacciones >= 20)

nrow(trabajos)
trabajos %>% arrange(n_transacciones) %>%
  select(archivo, n_transacciones, generar) %>%
  head(4)
sum(trabajos$generar)
\end{rcode}
\begin{rsalida}
[1] 40
# A tibble: 4 x 3
  archivo                          n_transacciones generar
  <chr>                                      <int> <lgl>  
1 2023-Q2_T07_sucursal-plaza-b.pdf              16 FALSE  
2 2023-Q4_T05_sucursal-oeste.pdf                18 FALSE  
3 2023-Q4_T07_sucursal-plaza-b.pdf              18 FALSE  
4 2023-Q4_T01_sucursal-centro.pdf               19 FALSE  
[1] 34
\end{rsalida}

Hay que generar 40 reportes (10 tiendas $\times$ 4 trimestres). La
combinación con menos transacciones es la Sucursal Plaza B en el segundo
trimestre (16). Seis combinaciones quedan por debajo de 20 y se marcan con
\codigo{generar = FALSE}, así que el bucle generaría 34 reportes; en el
script de automatización bastaría con recorrer
\codigo{filter(trabajos, generar)}.
\end{solucion}

\begin{solucion}{m8-7}
\codigo{split()} divide el detalle en una lista con un data frame por
tienda; los nombres de esa lista serán los nombres de las hojas. Con
\codigo{c()} ponemos la hoja \codigo{Resumen} al principio.

\begin{rcode}
library(tidyverse)
library(writexl)
library(readxl)

transacciones <- read_csv("datasets/transacciones.csv",
                          show_col_types = FALSE)
tiendas <- read_csv("datasets/tiendas.csv", show_col_types = FALSE)

detalle <- transacciones %>%
  left_join(tiendas %>% select(tienda_id, nombre_tienda),
            by = "tienda_id") %>%
  mutate(hoja = str_remove(nombre_tienda, "Sucursal ")) %>%
  arrange(tienda_id, fecha)

resumen <- detalle %>%
  group_by(hoja) %>%
  summarise(ventas = round(sum(total_transaccion), 2),
            transacciones = n(), .groups = "drop") %>%
  arrange(desc(ventas))

# Una hoja por tienda (sin la columna auxiliar "hoja")
por_tienda <- split(detalle %>% select(-hoja), detalle$hoja)

dir.create("resultados", showWarnings = FALSE)
archivo <- "resultados/ventas_por_tienda.xlsx"
write_xlsx(c(list(Resumen = resumen), por_tienda), path = archivo)

# Verificación
hojas <- excel_sheets(archivo)
length(hojas)
hojas
filas <- sapply(hojas[-1], function(h) nrow(read_excel(archivo, sheet = h)))
sum(filas)
\end{rcode}
\begin{rsalida}
[1] 11
 [1] "Resumen" "Centro"  "Este"    "Express" "Norte"   "Oeste"  
 [7] "Outlet"  "Plaza A" "Plaza B" "Premium" "Sur"    
[1] 1000
\end{rsalida}

El libro tiene 11 hojas (el resumen y diez tiendas) y las hojas de tiendas
suman las 1,000 transacciones: no se perdió ni se duplicó ninguna. Observa
que \codigo{split()} ordena las hojas alfabéticamente (Centro, Este,
Express\ldots); si quisieras otro orden, reordena la lista antes de
escribirla.
\end{solucion}

\begin{solucion}{m8-8}
La función valida sus argumentos con \codigo{stop()}, igual que la plantilla
del reporte ejecutivo. El bucle envuelve cada llamada en
\codigo{tryCatch()}: las filas exitosas se guardan en una lista de
resultados y los errores en el registro, así que ninguna combinación inválida
detiene el proceso.

\begin{rcode}
library(tidyverse)
library(scales)

transacciones <- read_csv("datasets/transacciones.csv",
                          show_col_types = FALSE)
tiendas <- read_csv("datasets/tiendas.csv", show_col_types = FALSE)

revisar_tienda <- function(id, mes) {
  if (!id %in% tiendas$tienda_id) stop("La tienda ", id, " no existe.")
  if (!mes %in% transacciones$mes) stop("No hay datos del mes ", mes)
  mes_previo <- format(as.Date(paste0(mes, "-01")) - 1, "%Y-%m")
  if (!mes_previo %in% transacciones$mes) {
    stop("No hay datos del mes anterior (", mes_previo, ")")
  }
  v <- transacciones %>% filter(tienda_id == id)
  actual   <- sum(v$total_transaccion[v$mes == mes])
  anterior <- sum(v$total_transaccion[v$mes == mes_previo])
  variacion <- actual / anterior - 1
  tibble(tienda_id = id,
         tienda = tiendas$nombre_tienda[tiendas$tienda_id == id],
         mes = mes, variacion = variacion,
         estado = case_when(variacion >  0.10 ~ "Crecimiento sólido",
                            variacion >= 0    ~ "Estable",
                            variacion > -0.10 ~ "Ligera caída",
                            TRUE              ~ "Alerta"))
}

pendientes <- tibble(id = c(tiendas$tienda_id, 11, 1),
                     mes = c(rep("2023-06", 10), "2023-06", "2023-01"))
resultados <- list()
registro <- list()

for (fila in seq_len(nrow(pendientes))) {
  id  <- pendientes$id[fila]
  mes <- pendientes$mes[fila]
  estado <- tryCatch({
    resultados[[length(resultados) + 1]] <- revisar_tienda(id, mes)
    "OK"
  }, error = function(e) paste("ERROR:", conditionMessage(e)))
  registro[[fila]] <- tibble(id = id, mes = mes, estado = estado)
}

registro <- bind_rows(registro)
registro %>% count(estado)
registro %>% filter(estado != "OK")

resultados <- bind_rows(resultados)
resultados %>%
  mutate(variacion = percent(variacion, accuracy = 0.1)) %>%
  select(tienda, variacion, estado)

# Texto Markdown para el reporte (iría en un bloque results = "asis")
alertas <- resultados %>% filter(estado == "Alerta")
cat(paste0("- **", alertas$tienda, "**: las ventas cayeron ",
           percent(abs(alertas$variacion), accuracy = 0.1),
           " respecto a mayo."), sep = "\n")
\end{rcode}
\begin{rsalida}
# A tibble: 3 x 2
  estado                                             n
  <chr>                                          <int>
1 ERROR: La tienda 11 no existe.                     1
2 ERROR: No hay datos del mes anterior (2022-12)     1
3 OK                                                10
# A tibble: 2 x 3
     id mes     estado                                        
  <dbl> <chr>   <chr>                                         
1    11 2023-06 ERROR: La tienda 11 no existe.                
2     1 2023-01 ERROR: No hay datos del mes anterior (2022-12)
# A tibble: 10 x 3
   tienda           variacion estado            
   <chr>            <chr>     <chr>             
 1 Sucursal Centro  -48.1%    Alerta            
 2 Sucursal Norte   -73.5%    Alerta            
 3 Sucursal Sur     31.8%     Crecimiento sólido
 4 Sucursal Este    245.5%    Crecimiento sólido
 5 Sucursal Oeste   -41.0%    Alerta            
 6 Sucursal Plaza A -16.3%    Alerta            
 7 Sucursal Plaza B 29.1%     Crecimiento sólido
 8 Sucursal Outlet  612.9%    Crecimiento sólido
 9 Sucursal Premium -8.0%     Ligera caída      
10 Sucursal Express -45.5%    Alerta            
- **Sucursal Centro**: las ventas cayeron 48.1% respecto a mayo.
- **Sucursal Norte**: las ventas cayeron 73.5% respecto a mayo.
- **Sucursal Oeste**: las ventas cayeron 41.0% respecto a mayo.
- **Sucursal Plaza A**: las ventas cayeron 16.3% respecto a mayo.
- **Sucursal Express**: las ventas cayeron 45.5% respecto a mayo.
\end{rsalida}

De las 12 combinaciones, 10 se revisaron y 2 quedaron en el registro con su
causa: la tienda 11 no existe y enero de 2023 no tiene mes anterior en los
datos. En junio, cinco tiendas están en alerta y la Sucursal Outlet creció
612.9\,\%: una variación tan grande con tan pocas transacciones merece
revisar los datos antes de celebrarla. El texto final está listo para
insertarse en un reporte con \codigo{results = "asis"}.
\end{solucion}
